\documentclass[aspectratio=169,10pt]{beamer}
\usetheme[numbering=fraction,progressbar=frametitle]{metropolis}
\usepackage{fontspec}
\setsansfont{Latin Modern Sans}
\usepackage{amsmath,amssymb,mathtools,graphicx,booktabs}
\usepackage{polyglossia}
\setdefaultlanguage{french}
\setbeamertemplate{navigation symbols}{}
\setbeamerfont{frametitle}{size=\large}
\setbeamerfont{normal text}{size=\normalsize}
\setlength{\parskip}{.45em}
\hypersetup{pdfauthor={},pdfsubject={Lecture slides},pdfkeywords={},pdftitle={Nombres complexes et géométrie}}
\providecommand{\N}{\mathbb{N}}
\providecommand{\Z}{\mathbb{Z}}
\providecommand{\D}{\mathbb{D}}
\providecommand{\Q}{\mathbb{Q}}
\providecommand{\R}{\mathbb{R}}
\providecommand{\C}{\mathbb{C}}
\providecommand{\vect}[1]{\overrightarrow{#1}}
\providecommand{\vi}{\vec{\imath}}
\providecommand{\vj}{\vec{\jmath}}
\providecommand{\vk}{\vec{k}}
\providecommand{\norm}[1]{\left\lVert #1\right\rVert}
\providecommand{\vc}[2]{\begin{pmatrix}#1\\#2\end{pmatrix}}
\providecommand{\deter}[2]{\det\left(#1,#2\right)}
\providecommand{\abs}[1]{\left\lvert #1\right\rvert}
\providecommand{\dd}{\,\mathrm{d}}
\providecommand{\e}{\mathrm{e}}
\providecommand{\Card}{\operatorname{Card}}
\providecommand{\E}{\operatorname{E}}
\providecommand{\V}{\operatorname{V}}
\providecommand{\degre}{^{\circ}}
\providecommand{\un}[1]{\,\mathrm{#1}}
\providecommand{\sci}[2]{#1\times10^{#2}}
\providecommand{\pgcd}{\operatorname{PGCD}}
\providecommand{\ppcm}{\operatorname{PPCM}}
\providecommand{\llbracket}{\mathopen{[\mkern-3mu[}}
\providecommand{\rrbracket}{\mathclose{]\mkern-3mu]}}
\title{Nombres complexes et géométrie}
\subtitle{Cours et exemples}
\date{}
\author{}
\begin{document}
\begin{frame}[plain]
\vfill
{\usebeamerfont{title}\inserttitle\par}
\medskip
{\usebeamerfont{subtitle}\insertsubtitle\par}
\medskip
{\usebeamercolor[fg]{progress bar}\rule{\textwidth}{.5pt}}
\vfill
\end{frame}
\begin{frame}{Objectifs}
Interpréter le module et un argument de \(\dfrac{c-a}{b-a}\) ; démontrer un alignement, une orthogonalité, la nature d'un triangle ; déterminer des ensembles de points ; déterminer les racines \(n\)-ièmes de l'unité et les utiliser pour les polygones réguliers.
\end{frame}
\begin{frame}{Activité d'introduction}
\begin{enumerate}
\setcounter{enumi}{0}\item Factoriser \(z^3-1\) par \(z-1\), puis résoudre \(z^3=1\) dans \(\C\). On note \(j\) la solution de partie imaginaire positive.

\item Placer les trois solutions dans le plan complexe. Calculer les distances entre elles. Quelle figure obtient-on ?

\item Écrire \(j\) sous forme exponentielle. Vérifier que \(j^2=\overline j\) et que \(1+j+j^2=0\). Interpréter cette dernière égalité à l'aide du centre de gravité.

\item Résoudre \(z^4=1\) et placer les solutions. Même question pour \(z^6=1\), en utilisant \(z^6-1=\left(z^3-1\right)\left(z^3+1\right)\).

\end{enumerate}
\end{frame}
\begin{frame}{Activité d'introduction (2)}
\begin{enumerate}
\setcounter{enumi}{4}\item Formuler une conjecture sur les solutions de \(z^n=1\) : leur nombre, leur module, leurs arguments, la figure qu'elles forment.

\end{enumerate}
\end{frame}
\begin{frame}{Angles et longueurs}
Soit \(A\), \(B\), \(C\) trois points d'affixes \(a\), \(b\), \(c\), avec \(A\neq B\) et \(A\neq C\). Alors
\[\abs{\dfrac{c-a}{b-a}}=\dfrac{AC}{AB}\qquad\text{et}\qquad\arg\left(\dfrac{c-a}{b-a}\right)\equiv\left(\vect{AB},\vect{AC}\right)\pmod{2\pi}.\]
\end{frame}
\begin{frame}{Angles et longueurs}
Le module d'un quotient est le quotient des modules : \(\abs{c-a}=AC\) et \(\abs{b-a}=AB\). Pour l'argument, \(\arg(b-a)\equiv\left(\vect u,\vect{AB}\right)\) et \(\arg(c-a)\equiv\left(\vect u,\vect{AC}\right)\), donc par la relation de Chasles \(\arg\dfrac{c-a}{b-a}\equiv\left(\vect u,\vect{AC}\right)-\left(\vect u,\vect{AB}\right)\equiv\left(\vect{AB},\vect{AC}\right)\pmod{2\pi}\).
\end{frame}
\begin{frame}{Angles et longueurs}
Avec les mêmes notations, posons \(q=\dfrac{c-a}{b-a}\).

\begin{enumerate}
\setcounter{enumi}{0}
\item \(A\), \(B\), \(C\) sont alignés si et seulement si \(q\) est réel.

\item Les droites \((AB)\) et \((AC)\) sont perpendiculaires si et seulement si \(q\) est imaginaire pur.

\item Le triangle \(ABC\) est rectangle isocèle en \(A\) si et seulement si \(q=i\) ou \(q=-i\) ; il est équilatéral si et seulement si \(q=\e^{i\pi/3}\) ou \(q=\e^{-i\pi/3}\).

\end{enumerate}
\end{frame}
\begin{frame}{Angles et longueurs}
\(q\) est réel si et seulement si son argument est \(0\) ou \(\pi\) modulo \(2\pi\), c'est-à-dire si \(\vect{AB}\) et \(\vect{AC}\) sont colinéaires. \(q\) est imaginaire pur (non nul) si et seulement si son argument est \(\pm\dfrac\pi2\). Pour 3, un triangle rectangle isocèle en \(A\) correspond à \(AC=AB\) et \(\left(\vect{AB},\vect{AC}\right)\equiv\pm\dfrac\pi2\) ; un triangle équilatéral à \(AC=AB\) et \(\left(\vect{AB},\vect{AC}\right)\equiv\pm\dfrac\pi3\).
\end{frame}
\begin{frame}{Angles et longueurs}
Soit \(A\) et \(B\) d'affixes \(a\) et \(b\). Le carré direct \(ABCD\) a pour sommets \(c=b+i(b-a)\) et \(d=a+i(b-a)\) : on passe de \(\vect{AB}\) à \(\vect{BC}\) par un quart de tour direct. Avec \(a=0\) et \(b=1\) : \(c=1+i\) et \(d=i\).
\end{frame}
\begin{frame}{Racines \(n\)-ièmes de l'unité}
Soit \(n\ge1\) un entier. On appelle racine \(n\)-ième de l'unité tout complexe \(z\) tel que \(z^n=1\). Leur ensemble est noté \(\mathbb U_n\).
\end{frame}
\begin{frame}{Racines \(n\)-ièmes de l'unité}
Pour tout entier \(n\ge1\), \(\mathbb U_n\) a exactement \(n\) éléments :
\[\mathbb U_n=\left\{\e^{2ik\pi/n}\;:\;k\in\llbracket 0;n-1\rrbracket\right\}.\]
\end{frame}
\begin{frame}{Racines \(n\)-ièmes de l'unité}
Soit \(z\in\mathbb U_n\). Alors \(\abs z^n=\abs{z^n}=1\) avec \(\abs z\ge0\), donc \(\abs z=1\) et \(z=\e^{i\theta}\) pour un réel \(\theta\). L'égalité \(\e^{in\theta}=1\) équivaut à \(n\theta\equiv0\pmod{2\pi}\), c'est-à-dire \(\theta=\dfrac{2k\pi}n\) avec \(k\) entier. Réciproquement, \(\left(\e^{2ik\pi/n}\right)^n=\e^{2ik\pi}=1\).
\end{frame}
\begin{frame}{Racines \(n\)-ièmes de l'unité (2)}
Deux entiers \(k\) et \(k'\) donnent le même nombre si et seulement si \(\dfrac{2(k-k')\pi}{n}\) est un multiple de \(2\pi\), soit \(k\equiv k'\pmod n\). Tout entier \(k\) a même reste qu'un entier de \(\llbracket 0;n-1\rrbracket\), et ces \(n\) entiers ont des restes distincts : \(\mathbb U_n\) a exactement \(n\) éléments.
\end{frame}
\begin{frame}{Racines \(n\)-ièmes de l'unité}
Pour \(n\ge3\), les images des éléments de \(\mathbb U_n\) sont les sommets d'un polygone régulier à \(n\) côtés inscrit dans le cercle unité, dont un sommet est le point d'affixe \(1\). Avec \(\omega=\e^{2i\pi/n}\), on a \(\mathbb U_n=\{1;\omega;\omega^2;\dots;\omega^{n-1}\}\). Cas à connaître : \(\mathbb U_2=\{1;-1\}\), \(\mathbb U_3=\{1;j;j^2\}\) avec \(j=\e^{2i\pi/3}\), \(\mathbb U_4=\{1;i;-1;-i\}\).
\end{frame}
\begin{frame}{Racines \(n\)-ièmes de l'unité}
Pour tout entier \(n\ge2\), la somme des racines \(n\)-ièmes de l'unité est nulle.
\end{frame}
\begin{frame}{Racines \(n\)-ièmes de l'unité}
\begin{columns}[T,onlytextwidth]
\begin{column}{.57\textwidth}
Avec \(\omega=\e^{2i\pi/n}\neq1\) (car \(n\ge2\)) : \(1+\omega+\dots+\omega^{n-1}=\dfrac{1-\omega^n}{1-\omega}=0\).
\end{column}
\begin{column}{.4\textwidth}
\centering
\includegraphics[width=\linewidth,height=4.5cm,keepaspectratio]{figures/expertes-05-1.pdf}
\end{column}
\end{columns}
\end{frame}
\end{document}
